\section{Usage}\label{sec:usage}

The pieces for using the checkpoints are published with this article in its results folder~\cite{ownfolder}: the reader in \UsageDirRead{}, the recipe of the render in \UsageDirRender{} and the converters of a map in \UsageDirConvert{}. This section says what each piece does and what was run and not run; the README of each folder says more.

Input. A checkpoint of Rectina reads a surface volume: a stack of layers sampled along the normal of a segmented surface, stored as an array of shape (planes, rows, columns), as an npy file or an OME-Zarr folder. The reader takes seventeen consecutive planes from the stack and slides a square patch over rows and columns. The stacks used in this work have twenty three planes, stored as uint16 holding the rendered grey value times 256, so that 65280 is the saturated value of an 8 bit render, and divided by 65280 before reading; the seven read starts are the seven windows of seventeen planes inside the twenty three. The reader applies to every patch the robust median absolute deviation normalisation of the villa ink detection code (percentiles 1 and 99). The stack has to be rendered along the surface normal at the scale of the training data, about 9 um. Reads of stacks from another render were not tested.

Reading a unit. The reader of \UsageDirRead{} is run with the villa package at the commit that its README pins (the commands of the README):

\begin{quote}
\begin{Verbatim}[fontsize=\scriptsize,breaklines=true,breakanywhere=true,breakautoindent=false]
git clone https://github.com/ScrollPrize/villa.git && git -C villa checkout 0bc2c241de7e289e36459468a9a0b32257679966
\end{Verbatim}
\end{quote}


\begin{quote}
\begin{Verbatim}[fontsize=\scriptsize,breaklines=true,breakanywhere=true,breakautoindent=false]
python3 rectina_read.py info    --checkpoint m2/ckpt_020000.pth
python3 rectina_read.py read    --stack unit.npy --checkpoint m2/ckpt_020000.pth --villa-src villa/vesuvius/src --device cuda --out out/unit
python3 rectina_read.py compare out/unit/map.npy reference_map.npy
\end{Verbatim}
\end{quote}


Each start reads its seventeen planes with two views (the identity and the view turned by 180 degrees, whose maps are averaged), four lattices shifted by a quarter of a patch, bfloat16 autocast on a graphics processor and Hann window blending; the map is the mean over the seven starts. The patch is \UsagePatchOne{} pixels for \UsageNameOne{} and \UsagePatchTwo{} pixels for \UsageNameTwo{}: it is the crop size of the configuration embedded in each checkpoint, read from the checkpoint (not inferred). A read needs a graphics processor in practice: for \UsageNameTwo{} a read of one million pixels took about \UsageSecPerMpxTwo{} seconds for the seven starts on a graphics processor of type \UsageDeviceTwo{} (peak memory \UsagePeakTwo{} MiB), and for \UsageNameOne{} about \UsageSecPerMpxOne{} seconds (peak memory \UsagePeakOne{} MiB), both read from the records of the reads of unit ag1. Run by: the author, on a rented RTX 4090 machine with bfloat16, for unit ag1 only, the commands read and compare of the README; the result is that the public reader reproduces the stored maps of the same checkpoints on this one unit (the sha256 of both maps equal the stored ones), and nothing more: it says nothing about other units, other stacks or other hardware. Also run: the comparison of a CPU read of a crop of unit ag1 with \UsageNameOne{} against the stored map, inside the declared limits. Not tested: the command info, \UsageNameTwo{} on a CPU, other graphics processors.

Output. For each unit the read writes map.npy (float32, rows by columns of the stack, values between 0 and 1), reads.csv and a record of the sha256 of the outputs. The map is a ranking score of ink against papyrus and not a calibrated probability: \UsageCalibration{} \LabelSmoothing{} On a real unit the values span a narrow band, for instance \UsageBandOne{} for \UsageNameTwo{} on ag1 and \UsageBandTwo{} on 0343, so a fixed threshold such as 0.5 means nothing; evaluate by rank (the area under the curve) or choose a threshold on a labelled region. To make an image, scale the map linearly between two percentiles and write an 8 bit tif.

Worked example. Unit ag1, a control unit of the comparison (organisers' human labels, none of its segments used in training), \UsageShapeOne{} pixels. The stored map of \UsageNameOne{} has sha256 \UsageShaOne{} and scores an area under the curve of \UsageAucOne{} (\UsageAucOneFull{}) against the labels with the scorer of this work; the stored map of \UsageNameTwo{} has sha256 \UsageShaTwo{} and scores \UsageAucTwo{} (\UsageAucTwoFull{}); both numbers are the cells of the comparison tables. Scoring a stored map needs no graphics processor. The declared limits for a map read on another device against a reference map are a largest absolute difference of \UsageTolAbs{}, a mean absolute difference of \UsageTolMean{} and a Pearson correlation of at least \UsageTolPearson{}.

Converting a map. The converter of \UsageDirConvert{} writes a map as a tif or as OME-Zarr and reads either back to check it against the npy (the commands of its README):

\begin{quote}
\begin{Verbatim}[fontsize=\scriptsize,breaklines=true,breakanywhere=true,breakautoindent=false]
python3 rectina_map_convert.py tif map.npy map_float32.tif
python3 rectina_map_convert.py tif map.npy map_uint16.tif --dtype uint16
python3 rectina_map_convert.py zarr map.npy map_v3.zarr
python3 rectina_map_convert.py zarr map.npy map_v2.zarr --zarr-format 2 --pixel-size-um 9.36
python3 rectina_map_convert.py check map.npy map_float32.tif
python3 rectina_map_convert.py check map.npy map_uint16.tif
python3 rectina_map_convert.py check map.npy map_v3.zarr
python3 rectina_map_convert.py check map.npy map_v2.zarr
\end{Verbatim}
\end{quote}


Run by: the test of that folder, on three real maps (two control units and one held-out unit): all four conversions and all four round trip checks passed. Not run: any OME-Zarr validator.

Rendering. The recipe of the render is in \UsageDirRender{}: the villa tool vc\_render\_tifxyz built from an upstream commit plus one patch of the study, with the settings of the stacks (twenty three planes at step one, float32 samples stored as uint16). It is a record of a recipe and does not replace the render component of the public chain. The binary of the original build is no longer on disk (a cleanup removed it); a rebuild from the same commit, patch and recipe has another file hash, and reproduced a stored render shard bit for bit in one reference run; the stacks of this recipe were not compared for equality with those of the public render component. Not tested: any other render of the same surface as input to the checkpoints.

Plugging in. The architecture is \texttt{\UsageArch{}} with a seventeen plane crop and \UsageTensors{} tensors. The ink detection package of villa builds it from the configuration embedded in the checkpoint (the function that makes the model, models/model.py of the package, which names this architecture). Verified by reading code only. Not tested: the inference command of that package with the checkpoints of this article, the current villa head.

What was not tested. Other renders, other scrolls than the fourteen units of the comparison, other graphics processors than the ones of the registered runs, meshes or surfaces that are not tifxyz, and scales other than the 9 um class of the training data.
